\documentclass[12pt]{article}
\usepackage[utf8]{inputenc}
\usepackage[T1]{fontenc}
\usepackage{hyperref}
\title{Parachi Language Resources for K-12 Curriculum Development}
\author{Compiled by Claude Sonnet 5.5 (claude-sonnet-5-5) for Andr\'as Kornai}
\date{}
\begin{document}
\maketitle
\section{Introduction}
Parachi is an Iranian language of the Parachi people in eastern Afghanistan (mainly upper Nijrab District), closely related to Ormuri. One source cites about 600 speakers and a UNESCO status of definitely endangered, which differs from the figure of about 3,500 in the request. It must not be confused with Pashai, a separate language.
\section{Primary Educational Resources}
None found.
\section{Linguistic and Reference Materials}
\begin{itemize}
\item Wikipedia article on Parachi (mentions V. A. Efimov's 2009 Russian grammar, texts and dictionary): \url{https://en.wikipedia.org/wiki/Parachi_language}
\item Episteme TLLM entry on Parachi: \url{https://episteme.tllm.fr/wiki/Parachi}
\end{itemize}
\section{Religious and Community Resources}
None found.
\section{Digital and Audio Resources}
None found.
\section{Curriculum Development Context}
The main scholarly description is Efimov's 2009 Russian grammar. No pedagogical materials or school programmes were found.
\end{document}